\section{Introduction}
\label{sec:intro}

\acresetall

GPU resources are a dominant cost of serving \acp{LLM}, motivating systems that improve accelerator utilization~\cite{orca,vllm,sarathi,distserve,splitwise,muxserve,aegaeon}. Latency-critical inference must meet stringent response-time \acp{SLO}, yet rarely uses every GPU resource at every instant.
Colocation systems reclaim this otherwise-idle capacity by running best-effort work, such as fine-tuning, alongside serving as long as its \acp{SLO} are preserved~\cite{orion,llmstation,harli}.

Many recent frontier \acp{LLM} use a \ac{MoE} architecture~\cite{deepseekv3,deepseekv4,qwen3,kimik2,gptoss,gemini15,gemini25}. An \ac{MoE} layer contains many feed-forward \emph{experts}, and a router sends each token to only a few of them~\cite{switch}, increasing model capacity without proportionally increasing per-token computation.
Since for large \acp{LLM} the total size of experts is beyond what a single GPU can store, such models are typically served with \ac{EP}, which places different experts on different GPUs~\cite{deepseekv3,vllm-ep-rfc}.
%and avoids sharding every expert across GPUs with \ac{TP}.


% \begin{figure}[t]
%   \centering
%   \inputplot{intro_motivation}
%   \caption{\textbf{Load balancing leaves GPUs idle.} Mean time per
%     \ac{MoE} layer on one GPU during a 128k-token prefill
%     (DeepSeek-V2-Lite, SGLang, 8$\times$H100, \ac{EP}$=8$). Even LPLB,
%     the strongest balancer we test, leaves each GPU waiting for 40\%
%     of the layer (\S\ref{sec:char-optimizations}).}
%   \label{fig:intro-motivation}
% \end{figure}

\ac{EP}, however, introduces a new source of underutilization.
Because token routing is input-dependent, GPUs receive different amounts of expert work.
Lightly loaded GPUs finish early and wait for the results of the most-loaded GPU before computation can proceed.
We call the time a GPU spends waiting for the most-loaded GPU a \emph{bubble}.
This load imbalance is well known and heavily optimized through techniques such as custom token routing objectives, strategic expert replication and placement, token redistribution~\cite{switch,gshard,expertchoice,deepseekv3,fastermoe,eplb,lplb,waterfill,harmoeny}, and specialized communication and expert kernels~\cite{deepep,pplxkernels,megablocks}.
We build on these mechanisms and ask the question they leave open:
\begin{highlight}
  \emph{Once load balancing and execution optimizations have done
  their job, how much idle GPU capacity remains, and can it be leveraged to perform useful work?}
\end{highlight}

% We call these stalls \emph{bubbles}.
% They arise mainly during \emph{prefill}, when the model processes the
% user's prompt, rather than during \emph{decode}, when it generates the
% answer one token at a time.
% Only prefill routes enough tokens per layer to turn imbalance into
% milliseconds of waiting.
% Bubbles grow with prompt length across the \ac{MoE} models we study,
% from the \textsc{Qwen} and \textsc{DeepSeek} families.
% On \textsc{Qwen3-30B-A3B}, the median bubble grows from \SI{0.22}{\milli\second}
% per layer at \SI{8}{k} tokens to \SI{1.2}{\milli\second} at \SI{32}{k} and
% \SI{6.4}{\milli\second} at \SI{128}{k}, adding up to \SI{59}{\milli\second} and
% \SI{307}{\milli\second} over its 48 \ac{MoE} layers.
%On \textsc{Qwen3-30B-A3B}, the median time spent in the bubble grows from 0.22\,ms per layer at
%8k tokens to 6.4\,ms at 128k, adding up to 307\,ms over its 48
%\ac{MoE} layers.
% Balancing beyond the strongest balancers we test
% (Figure~\ref{fig:intro-motivation}) needs more expert replicas, which take GPU memory
% from serving, or more token redistribution, which adds communication
% (\S\ref{sec:char-optimizations}).

% Fig. intro-pareto removed (user, 2026-10-23); kept here to restore if needed.
% \begin{figure}[t]
%   \centering
%   \inputplot{intro_pareto}
%   \caption{
%     Fine-tuning throughput versus serving latency, relative to serving
%     alone, for \sys{} and colocation baselines. Median over 45
%     configurations (three \ac{MoE} models, 4k--32k-token prompts, three
%     loads); bars show the interquartile range. \sys{} trains
%     2.6$\times$ faster than LLMStation at similar latency; naive
%     colocation either trains slowly or inflates latency.
%   }
%   \label{fig:intro-pareto}
%
% \end{figure}

\begin{figure}[t]
  \centering
  \inputplot{intro_trace}
  % \caption{
  %   Serving Qwen3-30B-A3B on a production trace (Azure LLM inference
  %   trace, code, original rate, \SI{30}{\second} bins) while fine-tuning
  %   it. Top: request arrivals. Bottom: fine-tuning throughput. Over the
  %   trace, \sys{} trains 3.3$\times$ faster on average than LLMStation and
  %   14$\times$ faster than \uncoord (40\% \ac{SM} cap), while
  %   its p90 \ac{TTFT} stays within 3\% of serving alone (legend:
  %   whole-trace p90 \ac{TTFT} relative to serving alone \oana{this paranthesis is very confusing. I would remove it, and instead annotate the lower plot, with different colors to mention the TTFT increase of x\% for each baseline and add to the legend something that explains that the TTFT is computed against serving alone. Right now, this important information is a bit hidden.}).%
  %   % \amk{here we could remove the TTFT plots and just leave with what is in the legend} \rnp{fixed, ptal! TTFT panel removed; whole-trace p90 TTFT is in the legend.}
  % }
  \caption{\textbf{\sys trains 3.3$\times$ faster than \llmstation and 14$\times$ faster than \uncoord on a production trace, while keeping p90 \ac{TTFT} within 3\% of serving alone.}
  Load (top) and fine-tuning throughput (bottom) of \sys and baselines over time on a \textsc{Qwen3-30B-A3B} model serving a production (Azure) trace at its original rate (with \SI{30}{\second} bins).
  Legend values give each system’s p90 \ac{TTFT} over the whole trace, relative to serving alone.}
  \label{fig:intro-trace}

\end{figure}

%\mv{We're not properly introducing colocation here.}
Our analysis in \S\ref{sec:characterization} answers the first half of this question and reveals that substantial idle capacity remains.
After applying state-of-the-art \ac{EP} load balancing techniques, bubbles still account for 8--13\% of prefill time.
The second half of the question, whether this capacity can perform useful work, is more challenging.
Existing colocation systems are natural candidates~\cite{flexllm,llmstation,harli,deltaserve}.
They place best-effort work, typically \ac{PEFT} such as \ac{LoRA}~\cite{lora}, on the same GPUs as inference.
%They place a best-effort job on the same GPUs as inference, typically \ac{PEFT} such as LoRA~\cite{lora}, which trains a small add-on to the model rather than the model itself and thus reuses the weights already in GPU memory.
However, they mostly exploit GPU underutilization during output generation (decode) and schedule work at the granularity of an inference or training step, whereas bubbles occur mostly during initial prompt processing (prefill), inside a single \ac{MoE} layer, and differ in length across GPUs.
Existing systems are therefore ineffective in leveraging such bubbles.
% They decide when to let it run so that inference keeps meeting its \acp{SLO}\rp{added a sentence}.
% However, they target a different kind of idle capacity. They mostly exploit the GPU left underused during decode, and decide what to run only between two consecutive passes of a batch through the
%model.
%\amk{Shrinked that, please check, thought it was too long for intro}\mv{Looks good to me}
% Their usual best-effort job, \ac{PEFT} such as \ac{LoRA}~\cite{lora},
% trains small adapters on top of the frozen model.
%However, these systems exploit a particular idle capacity, mostly GPU underutilization during decode, and operate at the granularity of a forward or backward pass.
%\amk{same here, check}\mv{Looks good!}
%However, bubbles mostly occur during prefill and within each \ac{MoE} layer. %, depending on the router's token assignment.
%Therefore, existing colocation systems are unable to exploit the idle capacity introduced by bubbles. % (\S\ref{sec:gap}).

%Bubbles, in contrast, open and close inside a single \ac{MoE} layer, and each GPU has a different one, depending on how many tokens the router just sent it. Running work at the right time is therefore not enough; it must also run on the right GPU. In a trace-driven experiment, a scheduler that knows the exact average bubble of each layer, but not which GPU owns it, recovers about 1\% of bubble time; deciding per GPU recovers 43--45\% (\S\ref{sec:gap}). Harvesting bubbles thus requires short units of work and a per-GPU decision made right after routing.


We introduce \sys, a colocation system that leverages bubbles to fine-tune \acp{LLM} with \ac{LoRA} while serving them.
The key insight behind \sys is that \emph{the length of a bubble is predictable}.
%\sys relies on bubbles being predictable.
%During inference, the router decides how many tokens each GPU receives before the experts run.
%During inference, the router decides how many tokens each GPU receives before the experts run, and with this count \sys predicts each GPU's bubble time with a median error of \SI{9}{\micro\second}, with the average bubble size being \SI{123}{\micro\second} (.
During inference, the router decides how many tokens each GPU receives before the experts run, and using this count, \sys predicts each GPU's bubble duration with a median error of \SI{8.9}{\micro\second}, about 7\% of the median bubble.
We introduce the first scheduler that admits fine-tuning work per GPU inside an \ac{MoE} layer, unlike existing colocation systems: \sys runs short chunks of fine-tuning work on each GPU only inside that GPU's predicted bubble window.
%Existing colocation systems, in contrast, decide what to run once per serving iteration, identically for all GPUs.
%\oana{I like how you mention the key insight but still I don't see what the novelty is in bubblehiker's technical contribution.} \rp{rewrote the previous two sentences, pls check}\mv{Looks good} %, which closes just before the bubble is predicted to end, thus minimizing the effect on end-to-e nd inference latency.
%Because windows differ across GPUs, each GPU's training work advances at its own pace.

We also design a compiler that turns the served model into training code.
A serving engine only runs the model's \emph{forward pass}, which pushes tokens through the layers to predict the next one.
Training also needs a \emph{backward pass}, which propagates the error back through the model to compute the gradients that update the \ac{LoRA} adapter.
The compiler takes the serving engine's forward pass, derives the matching backward pass, and splits the resulting training step into chunks short enough to fit in a single bubble.
The generated code operates directly on the weights already loaded for serving, so \sys fine-tunes the model without keeping a second copy of its weights in GPU memory.
% As a result, we do not need to manually write backward kernels or partition the training computation so that each piece fits in a bubble.
Existing colocation systems either train a second copy of the model~\cite{harli} or require porting each model architecture into their own training engine~\cite{flexllm,llmstation}, and they split training no finer than a layer or a serving iteration.
Our compiler instead works on the serving engine's model as-is, including its inference-only \ac{MoE} kernels and expert placement, so supporting a new architecture requires no hand-written kernels for backward passes.
%\oana{Same comment as paragraph above. Unclear what's new. Need to emphasize the compiler piece A LOT more, mentioning explicitly that the other engines need you to manually write backward passes and explaining why this is very annoying. Right now it reads a bit understated.} \rp{rewrote the first and last sentences, pls check; FlexLLM and LLMStation derive backward automatically, so we say ``porting'' rather than ``hand-written backward''}
%\svr{I think saying we don't need to manually partition the computation for bubble fit would be better than saying we don't need to hand-write the backward pass} \rp{rephrased, pls check}\mv{Looks good!}

% To fill bubbles, \sys must train the very model the GPUs are serving.
% But the serving engine can only run the model forward, and it does not compute
% gradients needed in training.
% A standard trainer could do that, but it loads a second copy of the model,
% doubling its memory footprint, and runs each training step in one go,
% whereas a bubble lasts only a few milliseconds.
% \sys therefore uses a compiler to turn the serving model into training
% code that works on the served weights and partitions the computation into chunks short
% enough to fit in a bubble.
% \amk{Not sure i understand this problem}\oana{It is unclear why you would need to hand-write the training implementation.}
% \svr{Reworked this, does it read better?}
% \mv{It's also somewhat unclear to me}\rp{rewrote, please check}
% %
% \TODO{Add final compiler LOC reduction and compiled-vs.-handwritten
% performance result once integrated measurements are complete.}
%\svr{I removed the comment about LoC reduction, and any discussions around hand-written vs compiler partitioning}

We implement \sys and compare it against two baselines.
%No existing inference--fine-tuning colocation system is designed for \ac{MoE} serving (\S\ref{sec:gap}).
% \mv{Parts of this is already discussed a few paragraphs before. There is too much emphasis here on explaining how the baseline works. The intro is already very long; maybe we could not discuss this in the intro?}\rp{shortened} \rnp{fixed, ptal! Rafael's shortening addressed ``too much emphasis on how the baseline works''; this edit addresses ``already discussed a few paragraphs before'': I removed the two remaining restatements of per-iteration planning (``these systems plan once per serving iteration, and none admits fine-tuning work inside an MoE layer'' and LLMStation ``schedules fine-tuning at serving-iteration boundaries''), since the colocation paragraph above already says it.}
Since there are no colocation systems specifically designed for \ac{MoE}, we port an existing colocation system, \llmstation~\cite{llmstation}, to support inference with \ac{MoE} models, and also compare with \uncoord, a stock trainer running beside serving without coordination.
%\rnp{I think this here is the easiest way to keep eval coherent with intro:
%Because harvesting bubbles requires a
%finer-grained training runtime than prior colocation systems provide,
%\sys{}'s end-to-end gains also include better use of conventional
%decode and idle-time headroom; \S\ref{sec:eval-ablation} separates
%this runtime benefit from the additional capacity provided by bubbles.}
%\rp{renamed naive colocation to uncoordinated colocation, macros \texttt{\textbackslash uncoord} and \texttt{\textbackslash Uncoord}}
% \amk{i would not talk about naive colocation, doesn't sound professional, may be colocation-only instead?} \oana{Need to explain that there are no systems specifically designed for collocating MoE. Otherwise, this reads as if you're creating baselines that you can beat.} \rnp{fixed, ptal!}
\Cref{fig:intro-trace} shows the fine-tuning throughput of \sys and both baselines while serving the \textsc{Qwen3-30B-A3B} \ac{LLM} on a bursty production trace.
Over the trace, \sys trains 3.3$\times$ faster than \llmstation and 14$\times$ faster than \uncoord, and it sustains this advantage throughout,
including during load bursts.
Meanwhile, the 90th-percentile \acf{TTFT} (p90 \ac{TTFT}) of \sys stays within 3\% of serving alone, whereas \uncoord increases it by 18\%.
Part of this gain comes from the fine-grained training runtime that bubbles require, which also makes better use of the decode and idle time that \llmstation already exploits.
Bubbles add the most at high load, when little decode and idle time is left for training, raising throughput by up to 27\% (\S\ref{sec:eval-ablation}).
%\rp{added these two sentences, pls check}\mv{Looks good to me}
Our further evaluation across three \ac{MoE} models and two production traces confirms that \sys turns idle time in \ac{MoE} serving into substantial fine-tuning throughput, while keeping both p90 \ac{TTFT} and \ac{TPOT} within 10\% of serving alone in 50 of 63 workloads.
%These results, and our other results using different \acp{LLM} and workloads, show that effectively \sys turns idle time in \ac{MoE} serving into substantial fine-tuning throughput, at almost no cost to inference latency.

% We measure tail inference latency, the 90th percentile (p90), both
% before the first answer token (\ac{TTFT}) and between subsequent
% tokens (\ac{TPOT}).
% Across three \ac{MoE} models and prompts of 1k--32k tokens, \sys{}
% trains 1.9--6.2$\times$ faster than LLMStation (median 2.8$\times$).
% It keeps both p90 \ac{TTFT} and \ac{TPOT} within 10\% of serving alone
% in 50 of 63 workloads, against 44 for LLMStation, with a median
% increase of 4\% and 6\%.
% Like LLMStation, \sys{} also uses decode and idle time.
% Bubbles add the most when serving load is high and that time is
% scarce, raising training throughput by 27--37\%
% (\S\ref{sec:eval-ablation}).
% On bursty production traces, \sys{} trains 3.3$\times$ faster than
% LLMStation while staying within 7\% of serving alone
% (Figure~\ref{fig:intro-trace}).
%

This paper makes the following contributions:
%\oana{If in need of space, can remove this list, as it repeats what we already said above}:
%\amk{Or keep only the bold text, because it is still nice to state clearly the contribution, intead we could in the tezt above says the first contribution is... second...etc}
%\rnp{Following Oana's and Annemarie's notes above: the bullets now only name each contribution and its section; the paragraphs above carry the details and numbers.}
\begin{itemize}
  \item \textbf{A characterization of expert-parallel bubbles}, showing that they survive
    load balancing, are predictable, and need per-GPU planning inside the \ac{MoE} layer (\S\ref{sec:characterization}).
  \item \sys, \textbf{a novel colocation system} that predicts each GPU’s bubble from the router’s decision, schedules short fine-tuning chunks only inside that window, and uses a compiler to derive those chunks automatically from the served \ac{LLM} (\S\ref{sec:design}).
  
  %\sys{}, which leverages them through router-aware prediction,
    %serving-controlled gated admission, and asymmetric execution of
    %compiler-derived training chunks (\S\ref{sec:design}). \rp{Too many compound terms here. The intro never explains ``gated admission'' or ``asymmetric execution'', so a reader cannot follow this bullet}\mv{Indeed, quite complicated, rephrased it.}
  \item \textbf{An evaluation of \sys and baselines} on three \ac{MoE} models and two production traces that breaks the gains down by source of capacity and by mechanism (\S\ref{sec:evaluation}).
\end{itemize}
